\mode*
\section{Debugging as self-directed learning}
\label{debugging-as-learning}

In this section we develop the theoretical argument: that debugging is
learning, that the debugger must generate the patterns of variation for
themselves, and that this ability is the same as the one exercised in the deep
approach to learning.

\subsection{A bug is an unknown}

A bug, from the programmer's point of view, is a discrepancy: the program
behaves in one way, the programmer expects another.
In variation-theoretic terms, the programmer has not discerned some critical
aspect of the program (or of the language, or of the environment)---if she
had, there would be no surprise.
Fixing the bug requires discerning exactly that aspect.
Debugging is thus learning: the object of learning is the not-yet-discerned
critical aspect, and the debugging session is over precisely when it has been
discerned.

\begin{frame}<presentation>
  \begin{idea}
    \begin{itemize}
      \item A bug is an unknown that we need to discover.
      \item The bug marks a critical aspect the programmer hasn't discerned.
      \item Debugging \(=\) learning that critical aspect.
    \end{itemize}
  \end{idea}
\end{frame}

But there is a crucial difference from the classroom: when debugging, there is
no teacher who has already identified the critical aspect and prepared a
pattern of variation for it.
The debugger must open up the dimensions of variation themselves.
Strictly speaking, there is one more option: the debugger can go find a
teacher---search for an explanation online, ask a colleague or, nowadays,
an AI assistant.
Then someone else provides the variation, or just the fix; we return to
what this means for the study in \cref{analysis,discussion}.

\subsection{Learning without a teacher: discovery}

\Textcite[Ch.~5]{NCOL} accounts for exactly this situation---learning where
nobody yet knows the answer---in his discussion of scientific discovery.
Science is collective learning, and the great discoveries can be described as
opening up new dimensions of variation.

\mode<presentation>{%
\begin{frame}
  \begin{example}<+>[Copernicus and science, {\cite[pp.~155--156]{NCOL}}]
    \begin{itemize}
      \item Science: collective learning.
      \item Copernicus' contribution was not the heliocentric view, it was
        \emph{centricity}.
      \item He opened the dimension of centricity, where heliocentric is a
        feature different from geocentric.
      \item But that allowed us to consider centricity of many other systems.
    \end{itemize}
  \end{example}

  \begin{example}<+>[Darwin and evolution, {\cite[pp.~157--160]{NCOL}}]
    \begin{itemize}
      \item Darwin also opened a dimension: species adapt, not replaced.
      \item Reproduced by Wallace doing the same thing independently.
    \end{itemize}
  \end{example}
\end{frame}%
}

Copernicus' contribution, in this account, was not the heliocentric view
itself but the dimension of \emph{centricity}: once that dimension is opened,
heliocentric and geocentric become two features that can vary, and the
centricity of other systems can be considered \autocite[pp.~155--156]{NCOL}.
Darwin similarly opened the dimension along which species adapt rather than
being replaced \autocite[pp.~157--160]{NCOL}.
The researcher is a learner who must introduce the variation themselves---and the scientific method, with its controlled experiments varying one thing
at a time, is a systematic way of doing precisely that.
We will see in \cref{debugging} that the generalisation pattern in particular
coincides with the scientific method, and that debugging follows the same
logic: the debugger runs experiments on the program, varying one aspect at a
time, until the critical aspect reveals itself.
The debugging literature draws this connection itself:
\textcite{WhyProgramsFail} teaches systematic debugging \emph{as} the
scientific method, creating and verifying hypotheses about the failure
through experiments.

\subsection{Learning without a teacher: the deep approach}

The second account of teacher-less learning in \citetitle{NCOL} is the deep
approach to learning (\cref{approaches}).
\Citeauthor{NCOL} argues that what distinguishes the deep approach is that the
learner attends to the object of learning---the meaning---and thereby
introduces the appropriate patterns of variation for themselves:
\textcquote[p.~145]{NCOL}{the learner has to create the necessary patterns of
variation and invariance. This is the deep approach to learning in terms of
the theory of variation}.
The claim rests on a series of studies that \citeauthor{NCOL} reviews
\autocite[pp.~142--151]{NCOL}: nursing students in a problem-based
programme took the initiative to seek variation in their learning
themselves---challenging their own perspective, seeking alternative
explanations and new angles, trying out different ways to understand
and to do things~\autocite[p.~265]{Silen2000} (the review presents
this as a study of medical students; the thesis itself studied the
nursing programme); the medical students of the follow-up study
likewise deliberately created variation in their own learning, down
to their reading~\autocite{FyreniusWirellSilen2007}; Chinese university
students' accounts of how they study exhibit the generalisation and contrast
patterns almost verbatim~\autocite{MartonWenWong2005}, while the students
exhibiting a surface approach instead \enquote{read it over and over again}---repetition without variation; and repeated readings of a difficult text
did not improve understanding unless the reader varied the focus and
perspective adopted~\autocite{MartonWenestam1988}.
% XXX(#1) Silen2000 (thesis full text) and FyreniusWirellSilen2007 (the
% author's doctoral kappa) are now read; MartonWenWong2005 and
% MartonWenestam1988 remain cited as reviewed in NCOL---no accessible
% copies found (see the search protocol); read them via the library
% before attributing anything to them directly.

Support for the idea that the ability to generate patterns of variation can be
acquired---and keeps working after instruction ends---comes from what
Marton calls generative learning, an observation due to
\textcite{HolmqvistGustavssonWernberg2007} across several learning studies.
As \citeauthor{NCOL} reports, \textcquote[pp.~227--228]{NCOL}{[t]he outcomes
of learning turned out to be better some time after the learning occasion than
immediately after it in certain classes}---those taught with patterns of
variation---with no such effect where patterns were not used
\autocite[pp.~227--228]{NCOL}; further demonstrations of the effect, across
several school subjects, are reviewed there as well
\autocite[p.~228]{NCOL}.
The students taught with patterns of variation apparently continued to
generate variation for themselves after the teaching ended.

\mode<presentation>{%
\begin{frame}
  \begin{block}{Generative learning~\autocite[pp.~227--228]{NCOL}}
    \begin{itemize}
      \item Some studies observed delayed learning effects.
      \item Outcomes better some time after instruction with patterns.
      \item No such effect for those instructed without patterns.
    \end{itemize}
  \end{block}
\end{frame}%
}

\subsection{Hypotheses}
\label{hypotheses}

Putting the pieces together: debugging is learning without a teacher; learning
requires patterns of variation; without a teacher the learner must generate
them; and the deep approach is exactly the disposition to generate them.
This yields our overarching hypothesis---that debugging skill and the deep
approach to learning are expressions of the same underlying ability---which
we operationalise as follows.

\begin{frame}
  \begin{restatable}{hypothesis}{hpatterns}\label{h:patterns}
    Successful debugging episodes exhibit the patterns of variation:
    contrast, generalisation and fusion.
    Unsuccessful or prolonged episodes lack one or more of the patterns.
  \end{restatable}
  \begin{restatable}{hypothesis}{hdeep}\label{h:deep}
    Students with a stronger deep approach generate more, and more
    complete, patterns of variation in their debugging episodes.
  \end{restatable}
  \begin{restatable}{hypothesis}{hsuccess}\label{h:success}
    Debugging success (fixing the bug, and doing so in fewer edit--run
    cycles) is positively associated with the deep approach and with the
    completeness of the generated patterns.
  \end{restatable}
\end{frame}

The hypotheses are our predicted answers to the research questions:
\cref{h:patterns} answers \cref{rq:patterns}---the patterns should appear,
and their presence should separate successful from unsuccessful
episodes---while \cref{h:deep} answers \cref{rq:deep} and \cref{h:success}
answers \cref{rq:success}.
\Cref{rq:experts} gets no hypothesis of its own: the expert comparison is
exploratory, although the overarching hypothesis suggests that the experts'
episodes should resemble the deep-approach students'---only more so.

A corollary, which follows from generative learning, is that teaching
programming with patterns of variation should over time also improve the
students' debugging---since debugging draws on the same ability to vary one
aspect while holding the others invariant.
We can probe this corollary within the course cohort: not all students use
the course material---some study on their own---so we can ask
\cref{rq:generative} again.

\begin{frame}
  \begin{restatable}{hypothesis}{hgenerative}\label{h:generative}
    Students who use the course material generate more, and more complete,
    patterns of variation in their debugging episodes than students who do
    not---controlling for approach and prior programming experience.
  \end{restatable}
\end{frame}

Thanks to generative learning, \cref{h:generative} predicts the answer to
\cref{rq:generative}.
Note that it follows from generative learning rather than from the
overarching hypothesis: \cref{h:deep,h:success} could hold while
\cref{h:generative} fails, and vice versa.
We return to this in \cref{discussion}.
